\chapter{Ekonomi Cloud, FinOps, dan Keberlanjutan}

\section*{Tujuan Pembelajaran}
\begin{itemize}
  \item Membandingkan struktur biaya pusat data sendiri dan \textit{cloud} dengan pendekatan TCO.
  \item Menjelaskan prinsip, fase, dan peran dalam FinOps Framework 2026.
  \item Merumuskan metrik biaya per unit dan metrik jejak karbon untuk satu layanan.
\end{itemize}

\section{Pendahuluan}

% TODO
% Cakupan: bagian keuangan Medikara Health terbiasa dengan anggaran server lima tahunan dan khawatir tagihan cloud bulanan tidak terkendali; bab ini menjelaskan cara memahami, memperkirakan, dan mengelola biaya cloud bersama-sama.

\section{CapEx, OpEx, dan Total Cost of Ownership}

% TODO
% Cakupan: perbedaan CapEx dan OpEx serta dampaknya pada laporan keuangan; komponen TCO, termasuk biaya yang sering terlupa (lisensi, tenaga kerja, migrasi, egress); cara menyusun perbandingan TCO yang jujur; business case disinggung singkat (dibahas rinci di IT40513).
% Kasus Medikara: perbandingan TCO kasar SIMRS satu rumah sakit di pusat data sendiri dan di cloud.

\section{Model Harga Cloud}

% TODO
% Cakupan: on-demand, komitmen pemakaian (reserved dan savings plan), dan kapasitas spot; risiko komitmen berlebih; pola beban yang cocok untuk tiap model harga.
% Pendalaman teknis: memakai kalkulator harga penyedia untuk membandingkan tiga model harga pada satu workload.

\section{FinOps Framework 2026}

% TODO
% Cakupan: definisi FinOps dan enam prinsipnya; fase Inform, Optimize, dan Operate; persona inti dan persona pendukung; empat domain dan capability, termasuk Executive Strategy Alignment yang baru pada 2026; scope dan kategori teknologi (public cloud, SaaS, AI, pusat data).
% Kasus Medikara: siapa di Medikara Health yang mengisi persona Finance, Engineering, Product, dan Leadership.

\section{Alokasi Biaya dan Unit Economics}

% TODO
% Cakupan: tagging sebagai dasar alokasi; showback dan chargeback; spesifikasi FOCUS 1.4 untuk data tagihan lintas penyedia; unit economics dan cara memilih unit yang bermakna bagi direksi.
% Kasus Medikara: biaya infrastruktur per kunjungan pasien dan per janji temu sebagai calon metrik.
% Pendalaman teknis: mengolah contoh data tagihan berformat FOCUS dengan spreadsheet atau SQL.

\section{Keberlanjutan dan Jejak Karbon}

% TODO
% Cakupan: sumber emisi dari workload digital; pilar keberlanjutan pada kerangka Well-Architected; SCI menurut ISO/IEC 21031:2024; pengaruh pilihan region, pemanfaatan sumber daya, dan jadwal kerja terhadap emisi; keterkaitan biaya dan karbon.

\section{Aktivitas: Model Biaya Awal dan Sasaran Nilai Cloud}

% TODO
% Tujuan aktivitas: menghasilkan perkiraan TCO kasar untuk satu workload, satu metrik biaya per unit, dan satu sasaran nilai cloud yang terukur.
% Peran dalam tim: anggota berlatar akuntansi atau keuangan memimpin struktur biaya dan asumsi; anggota teknis memperkirakan biaya sumber daya dengan kalkulator harga; anggota lain memeriksa kewajaran asumsi.
% Langkah: (1) pilih satu workload dari matriks Bab 2; (2) susun biaya kondisi sekarang; (3) perkirakan satu atau dua opsi cloud; (4) tetapkan metrik biaya per unit; (5) rumuskan sasaran nilai, misalnya biaya infrastruktur per kunjungan pasien turun 15 persen dalam dua tahun.
% Luaran dan tonggak capstone: C1, model biaya awal dan sasaran nilai cloud.
% Berkas sources yang direncanakan: sources/chapter03/templat-tco-sederhana.xlsx, sources/chapter03/contoh-tagihan-focus.csv, sources/chapter03/lembar-sasaran-nilai.md

\section{Ringkasan}

% TODO

\section{Latihan}

\begin{enumerate}
  \item Jelaskan mengapa tagihan \textit{cloud} sebuah organisasi bisa naik walaupun jumlah penggunanya tetap. Berikan tiga penyebab yang mungkin.
  \item Pilih satu layanan Medikara Health dan usulkan metrik biaya per unit yang bermakna bagi direksi. Jelaskan alasan pemilihannya.
  \item Kapan model harga berbasis komitmen lebih menguntungkan daripada \textit{on-demand}, dan kapan justru merugikan?
\end{enumerate}
